% ADR-003: Authentication Method
\subsection{ADR-003: OAuth2 via Google with JWT Tokens}

\begin{tabular}{p{3cm}p{10cm}}
\textbf{Status} & Accepted \\
\textbf{Date} & 2026-03-15 \\
\textbf{Decision} & Use Google OAuth2 for user authentication and JWT tokens for session management. \\
\end{tabular}

\textbf{Context}: The project requires token-based authentication with OAuth2 via Google as the primary provider. In a microservices architecture, session state must be shareable across services without a centralized session store bottleneck.

\textbf{Decision}: Implement the OAuth2 Authorization Code flow with Google as the identity provider. Upon successful authentication, the User Service issues a signed JWT containing the user's ID and roles. Subsequent requests include this JWT in the Authorization header. The API Gateway validates the JWT signature before routing requests to backend services.

\textbf{Token Strategy}:
\begin{itemize}
    \item Access Token: JWT with 15-minute expiry, signed with RS256.
    \item Refresh Token: Opaque token stored in Redis with 7-day expiry.
    \item Token refresh is handled by the User Service via the \texttt{/api/auth/refresh} endpoint.
\end{itemize}

\textbf{Consequences}:
\begin{itemize}
    \item \textit{Positive}: Stateless authentication at the API Gateway (no database lookup per request). Users benefit from Google's security infrastructure. JWTs are self-contained and can be verified by any service.
    \item \textit{Negative}: JWT revocation is not immediate (must wait for expiry or maintain a blacklist in Redis). Dependency on Google's OAuth2 service availability.
\end{itemize}
